%!TEX program = xelatex
%!TEX root = geometria_analitica.tex
%%Usar makeindex -s indexstyle.ist arquivo.idx no terminal para gerar o índice remissivo agrupado por inicial
%%Após executar pdflatex arquivo

\chapter{Vetores no Espaço}

Fixemos um ponto $O$ no espaço, que será denominado como \textbf{origem}. Tomemos três retas, duas a duas perpendiculares entre si e concorrentes em $O$,
que serão denominados \textbf{eixos coordenados} e serão denotados por $Ox$, $Oy$ e $Oz$.
\begin{figure}[!h]
    \centering
    \caption{Sistema de coordenadas tridimensionais}
    \includegraphics{eixos-coordenados-espaco.pdf}
\end{figure}

Projetando um ponto $P$ do espaço, ortogonalmente, sobre cada um dos eixos coordenados podemos representar os pontos do espaço por ternas ordenadas $(a,b,c)$
de números reais.

\begin{figure}[!h]
    \centering
    \caption{Coordenadas de um ponto no espaço tridimensional}
    \includegraphics{coordenadas-ponto-espaco.pdf}
\end{figure}

Assim os pontos de $\real^3$ são descritos por
\[
    \real^3 = \{(a,b,c)\mid a,b,c \in \real\}.
\]

Sejam $P(x_1,y_1,z_1)$ e $Q(x_2,y_2,z_2)$ dois pontos de $\real^3$. Traçando por $P$ um segmento $PS$ paralelo ao segmento $P_0Q_0$, onde $P_0(x_1,y_1)$ e
$Q_0(x_2,y_2)$, obtemos que

\[
    d(P,Q)^2 = d(P_0,Q_0)^2 + |SQ|^2.
\]

\begin{figure}
    \centering
    \caption{Distância entre pontos em $\real^3$}
    \includegraphics{distancia-pontos-espaco.pdf}
\end{figure}

Agora, $|SQ| = |z_2 - z_1|$ e daí
\[
    D(P,Q) = \sqrt{(x_1 - x_2)^2 + (y_1 - y_2)^2 + (z_1 - z_2)^2}.
\]

\begin{exemplos}
    A distância entre os pontos $P(2,-1,1)$ e $Q(-3,4,2)$ é
    \[
        d(P,Q) = \sqrt{[2 - (-3)]^2 + (-1-4)^2 + (1-2)^2} = \sqrt{51}.
    \]
\end{exemplos}

Os vetores e as operações com eles podem ser definidas utilizando-se dos eixos coordenados do espaço. Para isso, um ponto $P$ do espaço é representado em $\real^3$ por uma terna de coordenadas $P(x_1, y_1,z_1)$, onde $x_1$, $y_1$, $z_1 \in \real$.

Seja $\vec{u}$ um vetor no espaço. Como no caso do plano, podemos representar $\vec{u}$ como um orientado com origem no ponto $O(0,0,0)$ e extremidade em um ponto $(x_1,y_1,z_1)$.

Assim escrevemos
\[
    \vec{u} = (x_1,y_1,z_1).
\]
As coordenadas $(x_1,y_1,z_1)$ são chamadas de \textbf{componentes} do vetor $\vec{u}$. Com essa representação, o vetor nulo $\vec{0}$ é escrito como
\[
    \vec{0} = (0,0,0).
\]

Desse modo, podemos reescrever a definição de soma e produto por escalar da seguinte forma:
\begin{definicao}
    Sejam $\alpha \in \real$, $\vec{u} = (x_1, y_1,z_1)$ e $\vec{v} = (x_2, y_2,z_2)$. Então:
    \begin{enumerate}
        \item $\vec{u} + \vec{v} = (x_1 + x_2, y_1 + y_2,z_1 + z_2)$;
        \item $\alpha \vec{u} = (\alpha x_1, \alpha y_1, \alpha z_1)$.
    \end{enumerate}
\end{definicao}

Em alguns casos pode ser mais conveniente representar um vetor $\vec{u} = (x_1, y_1,z_1)$ de forma matricial. Para isso escrevemos as componentes de $\vec{u}$ como uma matriz de 1 coluna e três linhas
\[
    \vec{u} = \begin{bmatrix}
        x_1 \\y_1\\z_1
    \end{bmatrix}.
\]

Vimos que a norma de um vetor é definida como o comprimento de um segmento orientado que o represente. Com a representação de vetores por coordenadas podemos reescrever o conceito de norma da seguinte maneira:
\begin{definicao}
    Seja $\vec{u} = (x_1, y_1,z_1)$ um vetor. A \textbf{norma} de $\vec{u}$, denotada por $\norm{\vec{u}}$ é dada por
    \[
        \norm{\vec{u}} = \sqrt{x_1^2 + y_1^2 + z_1^2}.
    \]
\end{definicao}

Um vetor de norma 1 é chamado de \textbf{vetor unitário}. Dado um vetor não nulo $\vec{u}$ o vetor
\[
    \vec{v} = \dfrac{\vec{u}}{\norm{\vec{u}}}
\]
é um vetor unitário na direção de $\vec{u}$ pois
\[
    \norm{\vec{v}} = \norm{\dfrac{\vec{u}}{\norm{\vec{u}}}} = \dfrac{1}{\norm{\vec{u}}}\norm{\vec{u}} = 1.
\]

\begin{exemplo}
    Seja $\vec{u} = (-2,3,2)$ um vetor. Então
    \begin{align*}
        \norm{\vec{u}} & = \sqrt{(-2)^2 + 3^2 + 2^2} = \sqrt{17}                                                                                                      \\
        \vec{v}        & = \dfrac{\vec{u}}{\norm{u}} = \dfrac{1}{\sqrt{17}}(-2,3,2) = \left(\dfrac{-2}{\sqrt{17}}, \dfrac{3}{\sqrt{17}}, \dfrac{2}{\sqrt{17}}\right).
    \end{align*}
\end{exemplo}
% section vetores_plano (end)

\begin{proposicao}
    Sejam $\vec{u} = (x_1, y_1,z_1)$ um vetor em $\real^3$ e $\alpha \in \real$. Então:
    \begin{enumerate}
        \item Se $\vec{u} \ne \vec{0}$, então $\norm{\vec{u}} \ne 0$.
        \item $\norm{\vec{u}} \ne 0$ se, e somente se, $\vec{u} \ne \vec{0}$.
        \item $\norm{\alpha\vec{u}} = |\alpha|\norm{\vec{u}}$.
    \end{enumerate}
\end{proposicao}
\begin{prova}
    \begin{enumerate}
        \item  Se $\vec{u} \ne (0,0,0) = \vec{0}$, então $x_1 \ne 0$ ou $y_1 \ne 0$ ou $z_1 \ne 0$. Daí
              \[
                  \norm{\vec{u}} = \sqrt{x_1^2 + y_1^2 + z_1^2} > 0.
              \]
        \item $\norm{\vec{u}} = 0$ se, e somente se, $x_1^2 + y_1^2 + z_1^2 = 0$, isto é, $x_1 = y_1 = z_1 = 0$. Portanto $\vec{u} = \vec{0}$.
        \item $\norm{\alpha\vec{u}} = \norm{\alpha(x_1, y_1, \alpha z_1)} = \sqrt{(\alpha x_1)^2 + (\alpha y_1)^2 + (\alpha z_1)^2} = \sqrt{\alpha^2(x_1^2 + y_1^2 + z_1^2)} = |\alpha|\norm{\vec{u}}$.
    \end{enumerate}
\end{prova}

\subsection{Ângulo entre vetores e produto interno no espaço}

\begin{definicao}
    Sejam $\vec{u}$ e $\vec{v}$ vetores não nulos. O \textbf{ângulo} entre $\vec{u}$ e $\vec{v}$ é definido pelo ângulo $\theta$ determinado por $\vec{u}$ e $\vec{v}$ e que satisfaz $0 \le \theta \le \pi$, quando os vetores $\vec{u}$ e $\vec{v}$ são representados com a mesma origem.\index{Vetores!ângulo}
\end{definicao}

\begin{definicao}
    Quando o ângulo $\theta$ entre dois vetores $\vec{u}$ e $\vec{v}$ é reto, isto é, $\theta = \pi/2$, ou um deles é o vetor nulo, dizemos que os vetores $\vec{u}$ e $\vec{v}$ são \textbf{ortogonais} ou \textbf{perpendiculares} entre si. Denotamos tal fato, escrevendo $\vec{u} \perp \vec{v}$.\index{Vetores!Ortogonais}
\end{definicao}

Sejam $\vec{u}$ e $\vec{v}$ vetores não nulos e $\theta$ o ângulo entre eles.
\begin{figure}[!h]
    \centering
    \caption{ângulo entre vetores em $\real^3$}
    \includegraphics{angulo-vetores-espaco.pdf}
\end{figure}

\begin{definicao}\label{produtointerno-espaco}
    O \textbf{produto escalar} ou \textbf{produto interno} dos vetores $\vec{u}$ e $\vec{v}$, indicado por $\inner{u}{v}$, é o número real tal que\index{Vetores!Produto Escalar}
    \begin{enumerate}
        \item Se $\vec{u}$ ou $\vec{v}$ é nulo, então $\inner{u}{v} = 0$.
        \item Se $\vec{u}$ e $\vec{v}$ não são nulos e $\theta$ é o ângulo entre $\vec{u}$ e $\vec{v}$, então
              \begin{align}\label{produto-interno-espaco}
                  \inner{u}{v} = \norm{\vec{u}}\norm{\vec{v}}\cos\theta.
              \end{align}
    \end{enumerate}
\end{definicao}

\begin{proposicao}
    Sejam $\vec{u}$ e $\vec{v}$ vetores e $\theta$ o ângulo entre $\vec{u}$ e $\vec{v}$.
    \begin{enumerate}
        \item Se $\vec{u}$ e $\vec{v}$ não são nulos, então
              \[
                  \cos\theta = \dfrac{\inner{u}{v}}{\norm{\vec{u}}\norm{\vec{v}}}.
              \]
        \item Qualquer que seja o vetor $\vec{u}$,
              \[
                  \norm{\vec{u}} = \sqrt{\inner{u}{u}}.
              \]
        \item Quaisquer que sejam os vetores $\vec{u}$ e $\vec{v}$, $\vec{u}\perp\vec{v}$ se, e somente se, $\inner{u}{v} = 0$.
    \end{enumerate}
\end{proposicao}
De modo análogo ao caso para vetores no plano, temos:
\begin{teorema}
    O produto interno de dois vetores $\vec{u} = (x_1, y_1,z_1)$ e $\vec{v} = (x_2, y_2,z_2)$ em $\real^3$ é dado por
    \[
        \inner{u}{v} = x_1x_2 + y_1y_2 + z_1z_2.
    \]
\end{teorema}

\begin{exemplos}
    \begin{enumerate}
        \item Sejam $\vec{u} = (2, -3, 1)$ e $\vec{v} = (1,1, 0)$. Determine o ângulo entre $\vec{u}$ e $\vec{v}$.
              \begin{solucao}
                  Temos
                  \begin{align*}
                      \inner{u}{v} = \langle(2, -3, 1), (1,1,0)\rangle = 2 - 3 + 0 = -1 \\
                      \norm{\vec{u}} = \sqrt{14},\ \norm{\vec{v}} = \sqrt{2}            \\
                      \cos\theta = \frac{\inner{u}{v}}{\norm{\vec{u}}\norm{\vec{v}}} = -\dfrac{1}{\sqrt{14}\sqrt{2}} = -\dfrac{1}{\sqrt{28}}.
                  \end{align*}
                  Assim, $\theta = \arccos\left(-\dfrac{1}{\sqrt{28}}\right)$.
              \end{solucao}
        \item Encontre um vetor $\vec{u} = (x, y, z)$ ortogonal a $\vec{v} = (4, -1, 2)$ e tal que $\inner{u}{w} = -1$, onde $\vec{w} = (1,1,-1)$.
              \begin{solucao}
                  Como $\inner{u}{v} = 0$ e $\inner{u}{w} = -1$ obtemos o sistema
                  \[
                      \begin{cases}
                          4x - y + 2z = 0 \\
                          x + y - z = -1
                      \end{cases}
                  \]
                  cuja solução é $x = -1/5 - z/5$ e $y = -4/5 + 6z/5$. Fazendo $z = 0$, obtemos $\vec{u} = (-1/5, -4/5, 0)$.
              \end{solucao}
    \end{enumerate}
\end{exemplos}

\begin{proposicao}\label{propriedades-produto-vetorial-espaco}
    Quaisquer que sejam os vetores $\vec{u} = (x_1, y_1,z_1)$, $\vec{v} = (x_2, y_2,z_2)$ e $\vec{w} = (x_3, y_3,z_3)$ e qualquer que seja o número real $\lambda$ temos
    \begin{enumerate}[label=({\roman*})]
        \item\label{linearidade-soma-produto-vetorial-espaco} $\langle\vec{u}, \vec{v} + \vec{w}\rangle = \inner{u}{v} + \inner{u}{w}$
        \item\label{linearidade-escalar-produto-vetorial-espaco} $\langle\vec{u}, \lambda\vec{v}\rangle = \langle\lambda\vec{u}, \vec{v}\rangle = \lambda\inner{u}{v}$
        \item $\inner{u}{v} = \inner{v}{u}$
        \item Se $\vec{u} \ne \vec{0}$, então $\inner{u}{u} > 0$.
    \end{enumerate}
\end{proposicao}
\begin{prova}
    \begin{enumerate}[label=({\roman*})]
        \item \begin{align*}
                  \langle\vec{u}, (\vec{v} + \vec{w})\rangle & = \langle\vec{u}, (x_2 + x_3, y_2 + y_3, z_2 + z_3)\rangle = x_1(x_2 + x_3) + y_1(y_2 + y_3) + z_1(z_2+z_3) \\ &= (x_1x_2 + y_1y_2 + z_1z_2) + (x_1x_3 + y_1y_3 + z_1z_3) = \inner{u}{v} + \inner{u}{w}
              \end{align*}
        \item \begin{align*}
                  \langle\vec{u}, \lambda\vec{v}\rangle & = \langle\vec{u}, (\lambda x_2, \lambda y_2, \lambda z_2)\rangle = x_1(\lambda x_2) + y_1(\lambda y_2) + z_1(\lambda z_2) \\ &= (\lambda x_1)y_2 + (\lambda y_1)y_2 + (\lambda z_1)z_2 \\ &= \langle\lambda\vec{u}, \vec{v}\rangle\\
                  \langle\vec{u}, \lambda\vec{v}\rangle & = \langle\vec{u}, (\lambda x_2, \lambda y_2, \lambda z_2)\rangle = x_1(\lambda x_2) + y_1(\lambda y_2) + z_1(\lambda z_2) \\ &= \lambda (x_1y_2) + \lambda (y_1y_2) + \lambda(z_1z_2) \\ &= \lambda\inner{u}{v}
              \end{align*}
        \item $\inner{u}{v} = x_1x_2 + y_1y_2 + z_1z_2 = x_2x_1 + y_2y_1 + z_2z_1 = \inner{v}{u}$
        \item Se $\vec{u} \ne \vec{0}$, então $x_1 \ne 0$ ou $y_1 \ne 0$ ou $z_1 \ne 0$. Daí $\inner{u}{u} = x_1^2 + y_1^2 + z_1^2> 0$, como queríamos.
    \end{enumerate}
\end{prova}

\begin{observacao}
    \begin{enumerate}[label=({\alph*})]
        \item As propriedades \ref{linearidade-soma-produto-vetorial-espaco} e \ref{linearidade-escalar-produto-vetorial-espaco} da Proposição \ref{propriedades-produto-vetorial-espaco} podem ser estendidas para qualquer número de vetores:
              \[
                  \langle\vec{u}, \lambda_1\vec{v_1} + \lambda_2\vec{v_2} + \cdots + \lambda_n\vec{v_n}\rangle = \lambda_1\inner{u}{v_1} + \lambda_2\inner{u}{v_2} + \cdots + \lambda_n\inner{u}{v_n}.
              \]
        \item Na igualdade $\inner{u}{v} = \inner{u}{w}$ não podemos concluir que $\vec{v} = \vec{w}$. Por exemplo, para $\vec{u} = (1, 0, 0)$, $\vec{v} = (2, 1, 1)$ e $\vec{w} = (2, -5, 3)$ temos $\inner{u}{v} = \inner{u}{w}$ e no entanto $\vec{v} \ne \vec{w}$. Mas podemos concluir que $\vec{u}\perp(\vec{v} - \vec{w})$.
        \item De $\inner{u}{v} = 0$ não podemos concluir que $\vec{u} = \vec{0}$ ou $\vec{v} = \vec{0}$. Por exemplo, para $\vec{u} = (2, 1, 1)$ e $\vec{v} = (-2, 4, 0)$ temos $\inner{u}{v} = 0$ e no entanto $\vec{u} \ne \vec{0}$ e $\vec{v} \ne \vec{0}$.
    \end{enumerate}
\end{observacao}

\begin{proposicao}Sejam $\vec{u}$ e $\vec{v}$ vetores.
    \begin{enumerate}
        \item $\mid\inner{u}{v}\mid \le \norm{\vec{u}}\norm{\vec{v}}$ [Desigualdade de Scharwz]\index{Desigualdade de Scharwz}
        \item $\norm{\vec{u} + \vec{v}} \le \norm{\vec{u}} + \norm{\vec{v}}$ [Desigualdade Triangular]\index{Desigualdade Triangular}
    \end{enumerate}
\end{proposicao}
\begin{prova}
    Análoga à prova da Proposição \ref{DesigualdadeTriangular}.
\end{prova}

\subsection{Produto Vetorial}

Dados vetores $\vec{u} = (a_1,b_1,c_1)$ e $\vec{v} = (a_2,b_2,c_2)$, é possível encontrar um vetor $\vec{w} = (x,y,z)$ que seja simultaneamente ortogonal a $\vec{u}$ e a $\vec{v}$? Para isso devemos ter $\inner{u}{w} = 0$ e $\inner{v}{w} = 0$, ou seja,
\[
    \begin{cases}
        a_1x + b_1y + c_1z = 0 \\
        a_2x + b_2y + c_2z = 0
    \end{cases}.
\]
Podemos reescrever este último sistema como
\begin{equation}\label{sistemaprodutovetorial}
    \begin{cases}
        a_1x + b_1y = -c_1z \\
        a_2x + b_2y = -c_2z
    \end{cases}.
\end{equation}
Suponha que $\vec{u}$ e $\vec{v}$ não são paralelos, daí pelo menos um dos determinantes
\[
    d_{ab} = \det \begin{pmatrix}
        a_1 & b_1 \\
        a_2 & b_2
    \end{pmatrix}, \quad d_{ac} = \det \begin{pmatrix}
        a_1 & c_1 \\
        a_2 & c_2
    \end{pmatrix},\quad d_{bc} = \det \begin{pmatrix}
        b_1 & c_1 \\
        b_2 & c_2
    \end{pmatrix}
\]
é diferente de zero. Suponha então que $d_{ab} \ne 0$. Usando a Regra de Cramer, a solução do sistema \eqref{sistemaprodutovetorial} é dada por
\begin{align*}
    x = \dfrac{1}{d_{ab}}\det \begin{pmatrix}
                                  -c_1z & b_1 \\
                                  -c_2z & b_2
                              \end{pmatrix} = \dfrac{-z}{d_{ab}}\det \begin{pmatrix}
                                                                         c_1 & b_1 \\
                                                                         c_2 & b_2
                                                                     \end{pmatrix} = \dfrac{z(b_1c_2 - c_1b_2)}{d_{ab}}.
\end{align*}
Analogamente, obtemos
\[
    y = \dfrac{z(c_1a_2 - a_1c_2)}{d_{ab}}.
\]
Logo $\vec{w}$ é dado por
\[
    \vec{w} = \left(\dfrac{z(b_1c_2 - c_1b_2)}{d_{ab}}, \dfrac{z(c_1a_2 - a_1c_2)}{d_{ab}}, z\right),
\]
isto é, existem infinitos vetores $\vec{w}$ que são simultaneamente ortogonais a $\vec{u}$ e $\vec{v}$. Fazendo $z = d_{ab}$, obtemos que uma solução para o sistema \eqref{sistemaprodutovetorial} é dada por
\begin{equation}\label{produtovetorial}
    \vec{w} = (b_1c_2 - b_2c_1, a_2c_1 - a_1c_2, a_1b_2 - b_1a_2).
\end{equation}
O vetor $\vec{w}$ dado por \eqref{produtovetorial} é chamado de \textbf{produto vetorial} de $\vec{u}$ por $\vec{v}$ e é denotado por $\vec{u}\times\vec{v}$.\index{Vetores no Espaço!Produto Vetorial}

Um método alternativo para determinar o vetor $\vec{w}$ é o seguinte: denote por
\begin{align*}
    \vec{i} = (1,0,0) \\
    \vec{j} = (0,1,0) \\
    \vec{k} = (0,0,1).
\end{align*}
Então qualquer vetor $\vec{u} = (a_1,b_1,c_1)$ em $\real^3$ pode ser escrito como
\[
    \vec{u} = (a_1,b_1,c_1) = a_1(1,0,0) + b_1(0,1,0) + c_1(0,0,1) = a_1\vec{i} + b_1\vec{j} + c_1\vec{k}.
\]

Dados vetores $\vec{u} = (a_1,b_1,c_1)$ e $\vec{v} = (a_2,b_2,c_2)$, seja $\vec{w}$ o produto vetorial de $\vec{u}$ e $\vec{v}$. Sabemos que
\[
    \vec{w} = (b_1c_2 - b_2c_1, a_2c_1 - a_1c_2, a_1b_2 - b_1a_2).
\]

Agora, considere o determinante
\[
    \det \begin{bmatrix}
        \vec{i} & \vec{j} & \vec{k} \\
        a_1     & b_1     & c_1     \\
        a_2     & b_2     & c_2
    \end{bmatrix}.
\]
Temos
\begin{align*}
    \det \begin{bmatrix}
             \vec{i} & \vec{j} & \vec{k} \\
             a_1     & b_1     & c_1     \\
             a_2     & b_2     & c_2
         \end{bmatrix} & = b_1c_2\vec{i} + c_1a_2\vec{j} + a_1b_2\vec{k} - b_1a_2\vec{k} - c_1b_2\vec{i} - a_1c_2\vec{j} \\ &= (b_1c_2 - c_1b_2)\vec{i} + (c_1a_2 - c_2a_1)\vec{j} + (a_1b_2 - a_2b_1)\vec{k} \\ &= (b_1c_2 - c_1b_2, c_1a_2 - c_2a_1, a_1b_2 - a_2b_1) = \vec{w}.
\end{align*}

Logo
\[
    \vec{w} = \det \begin{bmatrix}
        \vec{i} & \vec{j} & \vec{k} \\
        a_1     & b_1     & c_1     \\
        a_2     & b_2     & c_2
    \end{bmatrix}
\]
onde $\vec{u} = (a_1,b_1,c_1)$ e $\vec{v} = (a_2,b_2,c_2)$.

\begin{exemplo}
    Sejam $\vec{u} = (-1,2,4)$ e $\vec{v} = (1,3,5)$. Determine $\vec{u}\times\vec{v}$ e $\vec{v}\times\vec{u}$.
    \begin{solucao}
        Temos
        \begin{align*}
            \vec{u}\times\vec{v} = \det \begin{bmatrix}
                                            \vec{i} & \vec{i} & \vec{k} \\
                                            -1      & 2       & 4       \\
                                            1       & 3       & 5
                                        \end{bmatrix} = -2\vec{i} + 9\vec{j} - 5\vec{k}.
        \end{align*}
        Logo $\vec{u}\times\vec{v} = (-2,9,-5)$.
        Agora,
        \begin{align*}
            \vec{v}\times\vec{u} = \det \begin{bmatrix}
                                            \vec{i} & \vec{i} & \vec{k} \\
                                            1       & 3       & 5       \\
                                            -1      & 2       & 4
                                        \end{bmatrix} = 2\vec{i} - 9\vec{j} + 5\vec{k}.
        \end{align*}
        Portanto $\vec{v}\times\vec{u} = (2,-9,5) = -\vec{u}\times\vec{v}$.
    \end{solucao}
\end{exemplo}

\begin{proposicao}
    Sejam $\vec{u}$, $\vec{v}$ e $\vec{w}$ vetores em $\real^3$ e $\alpha \in \real$. Temos
    \begin{enumerate}[label=({\alph*})]
        \item $\vec{u}\times\vec{v} = -\vec{v}\times\vec{u}$;
        \item $\vec{u}\times\vec{v} = \vec{0}$ se, e somente se, $\vec{u}\varparallel\vec{v}$;
        \item $\alpha(\vec{u}\times\vec{v}) = (\alpha\vec{u})\times\vec{v} = \vec{u}\times(\alpha\vec{v})$;
        \item $\vec{u}\times(\vec{v} + \vec{w}) = \vec{u}\times\vec{v} + \vec{u}\times\vec{w}$.
    \end{enumerate}
\end{proposicao}
\begin{prova}
    Seguem diretamente das propriedades do determinante.
\end{prova}

\begin{observacao}
    O sentido de $\vec{u}\times\vec{v}$ é dado pela regra da mão direita: se o ângulo entre $\vec{u}$ e $\vec{v}$ é $\theta$, giramos o vetor $\vec{u}$ de um ângulo $\theta$ até que coincida com o vetor $\vec{v}$ e acompanhamos o movimento com os dedos da mão direita, então o polegar vai apontar no sentido de $\vec{u}\times\vec{v}$.
\end{observacao}

\begin{proposicao}\label{normaprodutovetorial}
    Sejam $\vec{u}$ e $\vec{v}$ vetores no espaço. Então
    \begin{equation}
        \norm{\vec{u}\times\vec{v}}^2 = \norm{\vec{u}}^2\norm{\vec{v}}^2 - \inner{u}{v}^2.
    \end{equation}
\end{proposicao}
\begin{prova}
    Seja $\vec{u} = (a_1,b_1,c_1)$ e $\vec{v} = (a_2,b_2,c_2)$. Basta calcular $\norm{\vec{u}\times\vec{v}}^2$ e $\norm{\vec{u}}^2\norm{\vec{v}}^2 - \inner{u}{v}^2$  e comparar os termos.
\end{prova}

\begin{proposicao}
    Quaisquer que sejam os vetores não nulos $\vec{u}$ e $\vec{v}$ em $\real^3$, tem-se
    \[
        \norm{\vec{u}\times\vec{v}} = \norm{\vec{u}}\norm{\vec{v}}\sin\theta
    \]
    onde $\theta$ é o ângulo entre $\vec{u}$ e $\vec{v}$.
\end{proposicao}
\begin{prova}
    Da Proposição \ref{normaprodutovetorial}  e usando a equação \eqref{produto-interno-espaco} temos
    \begin{align*}
        \norm{\vec{u}\times\vec{v}}^2 & = \norm{\vec{u}}^2\norm{\vec{v}}^2 - \inner{u}{v}^2 \\ &= \norm{\vec{u}}^2\norm{\vec{v}}^2 - (\norm{\vec{u}}\norm{\vec{v}}\cos\theta)^2\\ &= \norm{\vec{u}}^2\norm{\vec{v}}^2(1 - \cos^2\theta)\\ &= \norm{\vec{u}}^2\norm{\vec{v}}^2\sin^2\theta.
    \end{align*}
    Portanto,
    \[
        \norm{\vec{u}\times\vec{v}} = \norm{\vec{u}}\norm{\vec{v}}\sin\theta
    \]
    pois $0 \le \theta \le \pi$.
\end{prova}

Agora considere o paralelogramo formado pelos vetores $\vec{u}$ e $\vec{v}$:
\begin{figure}[!h]
    \centering
    \caption{Interpretação geométrica do produto vetorial}
    \includegraphics{area-paralelogramo.pdf}
\end{figure}

A área $\mathcal{A}$ desse paralelogramo é dado Por
\[
    \mathcal{A} = bh.
\]
Neste caso $b = \norm{\vec{u}}$ e $h = \norm{\vec{v}}\sin\theta$. Logo
\[
    \mathcal{A} = \norm{\vec{u}}\norm{\vec{v}}\sin\theta = \norm{\vec{u}\times\vec{v}}.
\]
Portanto o vetor $\vec{u}\times\vec{v}$ é tal que seu módulo é numericamente igual \`a área do paralelogramo definido por $\vec{u}$ e $\vec{v}$.
\begin{exemplo}
    Calcule a área do triângulo de vértices $A(3,2,0)$, $B(0,4,3)$ e $C(1,0,2)$.
    \begin{solucao}
        Note que a área do triângulo é metade da área do paralelogramo. Assim seja
        \begin{align*}
            \vec{u} = \vec{CA} = (2,2,-2) \\
            \vec{v} = \vec{CB} = (-1,4,1)
        \end{align*}
        daí a área do triângulo será
        \[
            \mathcal{A} = \dfrac{\norm{\vec{u}\times\vec{v}}}{2}
        \]
        e como $\vec{u}\times\vec{v} = (10,0,10)$ temos $\mathcal{A} = 5\sqrt{2}$.
    \end{solucao}
\end{exemplo}

\subsection{Produto Misto}

Sejam $\vec{u}$ e $\vec{v}$ vetores de $\real^3$. Como $\vec{u}\times\vec{v}$ é um vetor de $\real^3$ podemos calcular seu produto interno com qualquer outro vetor $\vec{w}$ de $\real^3$. O número real
\begin{equation}
    \langle\vec{u}\times\vec{v}, \vec{w}\rangle
\end{equation}
é chamado de \textbf{produto misto} dos vetores $\vec{u}$, $\vec{v}$ e $\vec{w}$.\index{Vetores no Espaço!Produto Misto} Se $\vec{u} = (a_1,b_1,c_1)$, $\vec{v} =(a_2,b_2,c_2)$ e $\vec{w} = (a_3,b_3,c_3)$ temos
\begin{align*}
    \langle\vec{u}\times\vec{v}, \vec{w}\rangle & = \langle(b_1c_2 - c_2b_1,a_2c_1 - a_1c_2, a_1b_2 - a_2b_1), (a_3, b_3, c_3)\rangle \\ &= (b_1c_2 - b_2c_1)a_3 + (a_2c_1 - a_1c_2)b_2 + (a_1b_2 - a_2b_1)c_3.
\end{align*}

Agora,
\begin{align*}
    \det \begin{pmatrix}
             a_1 & b_1 & c_1 \\
             a_2 & b_2 & c_2 \\
             a_3 & b_3 & c_3
         \end{pmatrix} & = (-1)^{3+1}a_3\det \begin{pmatrix}
                                                 b_1 & c_1 \\
                                                 b_2 & c_2
                                             \end{pmatrix} \\ &+ (-1)^{3+2}b_3\det \begin{pmatrix}
        a_1 & c_1 \\
        a_2 & c_2
    \end{pmatrix} + (-1)^{3+3}c_3\det \begin{pmatrix}
        a_1 & b_1 \\
        a_2 & b_2
    \end{pmatrix}\\ &= a_3(b_1c_2 - c_1b_2) - b_3(a_1c_2 - a_2c_1) + c_3(a_1b_2 - a_2b_1).
\end{align*}

Portanto
\begin{equation}
    \langle\vec{u}\times\vec{v}, \vec{w}\rangle = \det \begin{pmatrix}
        a_1 & b_1 & c_1 \\
        a_2 & b_2 & c_2 \\
        a_3 & b_3 & c_3
    \end{pmatrix}.
\end{equation}

\begin{proposicao}
    Sejam $\vec{u}$, $\vec{v}$ e $\vec{w}$ vetores em $\real^3$. Então
    \begin{enumerate}[label=({\alph*})]
        \item $\langle\vec{u}\times\vec{v}, \vec{w}\rangle = \langle\vec{v}\times\vec{w}, \vec{u}\rangle$
        \item $\langle\vec{u}\times\vec{v}, \vec{w}\rangle = \langle\vec{u}, \vec{v}\times\vec{w}\rangle$
    \end{enumerate}
\end{proposicao}
\begin{prova}
    Segue diretamente das propriedades de determinantes.
\end{prova}

Sejam $\vec{u}$, $\vec{v}$ e $\vec{w}$ vetores em $\real^3$. Considere o paralelepípedo definido por $\vec{u}$, $\vec{v}$ e $\vec{w}$:
\begin{figure}[!h]
    \centering
    \caption{Interpretação geométrica do produto misto}
    \includegraphics{volume-paralelepipedo.pdf}
\end{figure}

A altura $h$ é dada por
\begin{align*}
    h = \norm{\vec{w}}|\cos\theta| = \norm{\vec{w}}\dfrac{|\langle\vec{u}\times\vec{v}, \vec{w}\rangle|}{\norm{\vec{u}\times\vec{v}}\norm{\vec{w}}} = \dfrac{|\langle\vec{u}\times\vec{v}, \vec{w}\rangle|}{\norm{\vec{u}\times\vec{v}}}.
\end{align*}
O volume do paralelepípedo é
\[
    \mathcal{V} = \norm{\vec{u}\times\vec{v}}h
\]
daí
\[
    \mathcal{V} = \norm{\vec{u}\times\vec{v}}\dfrac{|\langle\vec{u}\times\vec{v}, \vec{w}\rangle|}{\norm{\vec{u}\times\vec{v}}},
\]
isto é,
\begin{equation}
    \mathcal{V} = |\langle\vec{u}\times\vec{v}, \vec{w}\rangle|.
\end{equation}

\begin{exemplo}
    Os pontos $P(0,1,1)$, $Q(1,0,2)$, $R(1,-2,0)$ e $S(-2,2,-2)$ definem um paralelepípedo? Em caso afirmativo, qual o seu volume?
    \begin{solucao}
        Se os pontos $P$, $Q$, $R$ e $S$ definirem um paralelepípedo, seu volume deverá ser diferente de zero. Assim considere os vetores
        \begin{align*}
            \vec{u} = \vec{PQ} = (1,-1,1)  \\
            \vec{v} = \vec{PR} = (1,-3,-1) \\
            \vec{w} = \vec{PS} = (-2,1,-3).
        \end{align*}
        Assim temos
        \begin{align*}
            \langle\vec{u}\times\vec{v}, \vec{w}\rangle = \det \begin{bmatrix}
                                                                   1  & -1 & 1  \\
                                                                   1  & -3 & -1 \\
                                                                   -2 & 1  & -3
                                                               \end{bmatrix} = 0.
        \end{align*}
        Logo $P$, $Q$, $R$ e $S$ não definem um paralelepípedo.
    \end{solucao}
\end{exemplo}
